\subsection{アフィンLie作用の全重み量子化}
\label{subsec:tangent-quantization-v33}

整数 $m\ge3$ を取り、$x,y,z,w$ の重みを $(1,1,1,m)$ とする。
次の二つの導分を考える。
\[
 \mathcal H_m=\frac{(2m+3)x\partial_x+my\partial_y-3z\partial_z
                   -(4m+3)w\partial_w}{2(m+3)},
 \qquad \mathcal N_m=x\partial_y+2y\partial_z-z^m\partial_w.
\]
$[\mathcal H_m,\mathcal N_m]=\mathcal N_m/2$、$\operatorname{div}\mathcal H_m=-1/2$、
$\operatorname{div}\mathcal N_m=0$ であり、$\mathcal H_m\wedge \mathcal N_m$ は発散零のPoissonテンソルである。
$m=3$ では前節の例外成分の基準点のunimodular liftを回復する。
以下の量子化には、非可換な二つの導分に対する
Coll--Gerstenhaber--Giaquintoの普遍変形公式を用いる。
同じ公式の射影空間への適用はPym\cite[\S3.6, \S4]{Pym}に詳しい。

\begin{theorem}[アフィン作用からの全重みAS正則族]
\label{thm:tangent-star-v33}
$X=-2\mathcal H_m$、$N=\mathcal N_m$、$s=\hbar/2$ と置き、$A=\mathbb C[x,y,z,w]$ 上に
\[
 a\star b=\sum_{r\ge0}s^r N^r(a)\binom{X}{r}(b),
 \qquad
 \binom{X}{r}=\frac{X(X-1)\cdots(X-r+1)}{r!}
\]
を定める。この和は各 $a,b$ について有限であり、任意の
$\hbar\in\mathbb C$ に対して結合的な重み付き次数を保つ積を与える。
その通常の一次交換子は $\mathcal H_m\wedge \mathcal N_m$ である。

さらに $S_{\mathrm{aff},m}(\hbar)=(A,\star)$ はnoetherian整域かつ
global dimension四、GK次元四のAS正則代数である。
Hilbert級数とGorenstein parameterはそれぞれ
\[
 \frac1{(1-t)^3(1-t^m)},\qquad m+3
\]
である。特に $m=3$ の例外成分の基準点とその重み付き座標変換から得られる族は、
実際の非可換重み付き射影空間の座標環による量子化を持つ。
\end{theorem}

\begin{proof}
$[N,X]=N$ なので、上の普遍変形公式により結合性が従う。
補助次数を $\deg'(x,y,z,w)=(0,1,2,2m+1)$ とすると、$N$ は次数を一つ下げる。
従って $N$ は局所冪零であり、和は有限になる。
また $X,N$ はどちらも元の重み $(1,1,1,m)$ に関して次数零である。
一次項は $sN(a)X(b)=-\hbar N(a)\mathcal H_m(b)$ なので、
\[
 (a\star b-b\star a)/\hbar\bmod\hbar
 =\mathcal H_m(a)N(b)-N(a)\mathcal H_m(b).
\]
この計算で半古典括弧の符号を固定する。

AS正則性のために、星積が反復Ore拡大であることを示す。
$T=1+sN$ と置くと、$[N,X]=N$ と二項演算子の恒等式から
\[
 N(a\star b)=N(a)\star b+a\star N(b)+sN(a)\star N(b).
\]
従って $T$ は星積の自己同型である。局所冪零性により
$T^r=\sum_{k\ge0}\binom r k s^kN^k$ は任意のスカラー $r$ に対して定義できる。
$T^r(a\star b)=T^r(a)\star T^r(b)$ は非負整数 $r$ で成立し、
各 $a,b$ を固定すると両辺は $r$ の多項式なので、全ての $r$ で成立する。

変数を $x,y,z,w$ の順に付け加える。
この順で各 $N(t)$ は先行変数の多項式であり、
\[
 X(t)=\chi_t t,
 \qquad(\chi_x,\chi_y,\chi_z,\chi_w)
 =\frac1{m+3}\left(-(2m+3),-m,3,4m+3\right).
\]
先行変数の星積部分代数を $A_{<t}$ と書けば、
\[
 a\star t=T^{\chi_t}(a)\,t\qquad(a\in A_{<t})
\]
であり、右辺の積は通常の可換積である。
$\tau_t=T^{-\chi_t}|_{A_{<t}}$ と取ると
\[
 \partial_t(a):=t\star a-\tau_t(a)\star t\in A_{<t}.
\]
これは $N^k(t)\in A_{<t}$ ($k\ge1$) と
$\tau_t(a)\star t=at$ から従う。
結合性により $\partial_t$ は $\tau_t$-導分である。
また $t$ の星積冪は通常の冪 $t^j$ に関してmonicな三角変換をなし、
左係数には $T^{j\chi_t}$ の可逆な変換が作用する。
従って先行変数を係数とした左Ore基底が得られる。
これにより星積代数は、$\mathbb C[x]$ から三回の自己同型による
Ore拡大を行った代数と同型である。

Ore拡大はnoetherian性と整域性を保ち、
Liu--Wang--Wu\cite[Theorem 3.3 and Lemma 1.3]{LWW}によって
連結次数付きのAS正則代数となり、そのglobal dimensionは四である。
基礎の次数付きベクトル空間は可換多項式環と同じなので、Hilbert級数は表示通りであり、
GK次元は四である。Hilbert級数の相反性からGorenstein parameterは $m+3$ となる。
\end{proof}

\begin{corollary}[標準例外列と四対象の許容成分]
\label{cor:tangent-six-v33}
$S=S_{\mathrm{aff},m}(\hbar)$ に対し、
$(\mathcal O_S,\ldots,\mathcal O_S(m+2))$ は
$D^b(\operatorname{qgr}S)$ のfull strong exceptional collectionである。
従って $\Perf(S_{(0,1,m,m+1)})$ はその許容部分圏である。
$m=3$ では、元のVK四項列と標準六項列による実現を得る。
\end{corollary}

\begin{proof}
前定理とOrlov\cite[Corollary 2.7]{Orlov2009}による。
\end{proof}

ここで証明した量子化はAS正則であり、Ore拡大定理からtwisted Calabi--Yauでもある。
そのNakayama自己同型が恒等かどうかは本稿では決定していない。
したがって、この族の4-Calabi--Yau性は主張しない。
発散零の一次Poissonテンソルを持つことだけから、与えられた星積の
高次のmodular classやNakayama自己同型が消えるとは結論できない。

$m>3$ における上のPoissonテンソルが新たな既約成分を生成するかどうかは、
この量子化の定理からは結論しない。
